\subsection{智能体工作流设计}

\subsubsection{工作流执行过程}

智能体协作工作流的执行过程是一个多轮迭代的"分类-验证-修复"循环，其整体流程如图~\ref{fig:workflow-execution}所示。工作流从 Classifier Agent 开始，首先对目标函数进行语义分类，确定函数类别和推荐的替换策略。对于需要替换的函数，Classifier Agent 生成替换代码后，会调用 VerifyReplacement 进行代码审查和编译验证。

\begin{figure}[htbp]
\centering
\begin{tikzpicture}[
    node distance=1.5cm,
    box/.style={rectangle, draw, rounded corners, minimum width=2.5cm, minimum height=0.8cm, align=center, font=\small},
    decision/.style={diamond, draw, aspect=2, minimum width=2cm, align=center, font=\small},
    arrow/.style={->, thick, >=stealth}
]
    \node[box, fill=blue!20] (start) {开始：\\Classifier Agent};
    \node[box, fill=green!20, right of=start, xshift=0.5cm] (classify) {函数分类与\\替换代码生成};
    \node[decision, fill=yellow!20, right of=classify, xshift=0.8cm] (has_replacement) {需要\\替换?};
    \node[box, fill=orange!20, right of=has_replacement, xshift=0.8cm] (verify) {VerifyReplacement\\验证};
    \node[decision, fill=yellow!20, right of=verify, xshift=0.8cm] (pass) {验证\\通过?};
    \node[box, fill=red!20, above right of=pass, xshift=0.4cm] (fixer) {BuildFixer\\修复};
    \node[decision, fill=yellow!20, below right of=pass, xshift=0.4cm] (fix) {修复\\成功?};
    \node[box, fill=gray!20, right of=fixer, xshift=0.8cm] (accept) {采纳替换\\代码};
    \node[box, fill=gray!20, right of=pass, xshift=0.8cm] (end) {完成};
    
    \draw[arrow] (start) -- (classify);
    \draw[arrow] (classify) -- (has_replacement);
    \draw[arrow] (has_replacement) -- node[above, font=\scriptsize] {否} (end);
    \draw[arrow] (has_replacement) -- node[below, font=\scriptsize] {是} (verify);
    \draw[arrow] (verify) -- (pass);
    \draw[arrow] (pass) -- node[above, font=\scriptsize] {是} (accept);
    \draw[arrow] (pass) -- node[below, font=\scriptsize] {否} (fixer);
    \draw[arrow] (fixer) -- (fix);
    \draw[arrow] (fix) -- node[above, font=\scriptsize] {是} (accept);
    \draw[arrow] (fix) -- node[below, font=\scriptsize] {否} (end);
    \draw[arrow] (accept) -- (end);
\end{tikzpicture}
\caption{智能体协作工作流执行流程图}
\label{fig:workflow-execution}
\end{figure}

\subsubsection{智能体角色与职责}

替换代码生成后，系统需要进行质量验证，确保生成的代码能够在编译和运行阶段正常工作。本研究设计了两个层次的验证机制：代码审查和编译测试，形成"轻量检查→完整验证"的渐进式质量保障体系。

\paragraph{代码审查}

代码审查在代码生成后立即进行，目的是在编译之前就过滤掉明显有问题的替换代码，避免浪费编译时间和资源。

本研究采用 LLM 驱动的代码审查方法。该方法的核心思想是：将原函数代码和替换代码一起提交给 LLM，让 LLM 扮演代码审查者的角色，根据预定义的规则进行判断。LLM 具备理解代码语义的能力，能够识别替换代码是否保留了原函数的核心行为，是否引入了潜在的问题。

代码审查的规则通过提示词（Prompt）定义，主要包括以下方面：

\textbf{签名一致性规则}。要求替换函数的名称、返回类型、参数列表必须与原函数完全一致，这是替换代码能够正确集成回工程的基本前提。

\textbf{语义保持性规则}。要求替换代码应保留原函数的核心业务逻辑，如状态变量更新、错误处理流程等，只替换与硬件交互的部分。

\textbf{代码合理性规则}。要求替换代码不应引入明显的问题，如未声明变量、无限循环、内存泄漏等。

\textbf{调用关系保持规则}。如果原函数调用了某些核心函数（如回调函数、状态更新函数），替换代码应在适当的位置保留这些调用。

代码审查的实现采用 LLM 结构化输出技术。系统定义了 \texttt{CodeReviewResult} 数据类作为输出格式，包含 \texttt{passed}（是否通过检查）和 \texttt{reason}（未通过的原因）两个字段。通过 LangChain 的结构化输出功能，强制 LLM 按照该格式输出，确保结果可以被程序正确解析。

\paragraph{编译测试}

代码审查通过后，系统进行编译测试。编译测试是替换代码质量验证的关键环节，它将替换代码集成回工程，调用构建系统执行完整的编译链接流程，验证替换代码的语法正确性和链接兼容性。

编译测试的流程包括四个阶段：

\textbf{第一阶段：代码集成}。系统将替换代码写入源文件，替换原函数的定义。

\textbf{第二阶段：构建执行}。系统调用项目的构建脚本（如 CMake、Makefile）执行编译。

\textbf{第三阶段：错误解析}。如果编译失败，系统需要解析编译器的错误输出，提取错误类型、错误位置和错误原因。

\textbf{第四阶段：错误修复}。系统根据错误信息触发 Fixer Agent 进行自动修复，修复后重新编译，直到成功或达到重试上限。

编译集成是驱动函数分析、替换与编译模块的最后一个环节，负责将替换后的代码集成回工程并执行编译。这一环节由 Builder Agent 负责，它是系统中专门处理构建任务的智能体。

Builder Agent 采用 LangGraph 框架实现，具有状态管理和循环执行的能力。Agent 的工作流程是一个迭代过程：首先尝试编译，如果失败则分析错误并尝试修复，修复后重新编译，直到成功或达到重试上限。

Builder Agent 的核心能力包括：

\textbf{（1）构建环境配置}。读取项目的构建配置（如 CMakeLists.txt、Makefile），了解项目的构建方式和依赖关系。

\textbf{（2）编译错误分析}。当编译失败时，从编译器的错误输出中提取关键信息，判断错误的类型和严重程度。

\textbf{（3）修复触发}。当编译失败时，触发 Fixer Agent 进行错误修复。

\textbf{（4）产物验证}。编译成功后验证 ELF 文件的完整性。

\begin{figure}[htbp]
\centering
\begin{tikzpicture}[
    box/.style={rectangle, draw, rounded corners, minimum width=2.5cm, minimum height=1cm, align=center, font=\small},
    arrow/.style={->, thick, >=stealth}
]
    \node[box, fill=blue!20] (compile) at (0, 0) {编译\\执行};
    \node[box, fill=orange!20] (parse) at (4, 0) {错误\\解析};
    \node[box, fill=green!20] (fix) at (8, 0) {自动\\修复};
    
    \draw[arrow] (compile) -- (parse) node[midway, above, font=\scriptsize] {失败};
    \draw[arrow] (parse) -- (fix) node[midway, above, font=\scriptsize] {定位};
    \draw[arrow] (fix) -- (compile) node[midway, below, font=\scriptsize] {重试};
\end{tikzpicture}
\caption{编译错误修复闭环}
\label{fig:build-loop}
\end{figure}

\subsubsection{LangGraph 工作流框架}

系统采用 LangGraph 框架构建图式工作流。LangGraph 是 LangChain 生态中用于构建有状态、可循环的智能体工作流的库，支持复杂的状态管理和条件分支。

智能体工作流的核心包含三个逻辑节点：\textbf{推理与决策节点}（由大语言模型驱动）、\textbf{工具调用节点}（当推理节点需要补充信息时请求调用工具）、\textbf{结果汇总节点}（对输出进行结构化解析并写入缓存）。

\begin{figure}[htbp]
\centering
\begin{tikzpicture}[
    box/.style={rectangle, draw, rounded corners, minimum width=3cm, minimum height=1.5cm, align=center, font=\small},
    arrow/.style={->, thick, >=stealth}
]
    \node[box, fill=blue!20] (reason) at (5, 2) {推理与决策\\节点};
    \node[box, fill=green!20] (tool) at (0, 0) {工具调用\\节点};
    \node[box, fill=orange!20] (summary) at (10, 0) {结果汇总\\节点};
    
    \draw[arrow] (reason.south west) -- (tool.north east) node[midway, above, font=\scriptsize, sloped] {需要工具};
    \draw[arrow] (tool.north) -- (reason.south) node[midway, left, font=\scriptsize] {返回结果};
    \draw[arrow] (reason.south east) -- (summary.north west) node[midway, above, font=\scriptsize, sloped] {给出答案};
\end{tikzpicture}
\caption{智能体工作流架构图}
\label{fig:agent-workflow}
\end{figure}

\subsubsection{状态机设计}

LangGraph 的核心是状态（State）对象，它在工作流的各个节点之间传递和更新。状态模型需要维护智能体执行过程中的关键信息。

状态转移逻辑控制智能体的执行流程。智能体从 START 状态开始，首先进入推理节点，LLM 根据当前状态决定下一步行动。如果 LLM 请求调用工具（如查询函数源码、获取调用关系），则转移到工具节点，工具执行完毕后返回推理节点继续处理。如果 LLM 认为已经获得足够信息并给出最终答案，则转移到结果节点，对输出进行结构化解析后进入 END 状态。如果在执行过程中发生错误，且重试次数未超过上限，则返回推理节点重新尝试；否则进入错误处理流程。

\begin{figure}[htbp]
\centering
\fbox{\parbox{0.92\linewidth}{\centering
智能体状态机转移图（示意）
}}
\caption{智能体状态机转移图}
\label{fig:agent-state-machine}
\end{figure}
